\section{Related Work}
\label{sec:related_work}

\subsection{Language-Conditioned Visuomotor Policies}

Most language-conditioned visuomotor policies employ transformers to process
2D visual inputs and directly generate 3D actions for manipulation
\cite{brohan2022rt,brohan2023rt,kimopenvla,black2024pi_0,
li2023vision}.
Among these approaches, developing large vision-language-action (VLA)
models, most often by leveraging pre-trained vision-language models (VLMs),
has become increasingly popular because of their effectiveness in learning
complex manipulation skills
\cite{brohan2023rt,kimopenvla,li2023vision,black2024pi_0,
pertsch2025fast,intelligence2025pi_06star,intelligence2026pi_07,
generalist2025gen0,generalist2026gen1,genesis2026gene265,
sunday2026act2preview}.
However, such 2D image-based policies typically require substantial
data-collection effort, often relying on large trajectory datasets to
generalize effectively across tasks~\cite{o2024open}.
In contrast, 3D manipulation policies have demonstrated strong potential for
data-efficient learning by exploiting the spatial structure inherent in 3D
observations.
One line of work directly processes point clouds
\cite{chen2023polarnet,yuan2023m2t2,yang2025fp3,gervet2023act3d,
3d-da,garcia2024towards}.
For example, Act3D~\cite{gervet2023act3d} constructs a 3D feature cloud by
lifting image features onto the observed point cloud and predicts
translational actions by classifying candidate 3D points in the workspace.
Another line of work represents the observation space using voxels and
predicts translational actions within the same voxel space, thereby aligning
input observations and output actions in a shared spatial representation
\cite{shridhar2023perceiver,james2022coarse}.
More recently, RVT~\cite{goyal2023rvt} and
RVT-2~\cite{goyal2024rvt} leverage orthographic projections of 3D point
clouds to convert 3D signals into 2D images, avoiding the high computational
cost of directly processing native 3D representations.
Unlike the above methods, our base framework, \method, seeks to unify the
semantic effectiveness of VLA models with the data efficiency of 3D
manipulation policies within a single cohesive framework.


\subsection{3D Vision-Language-Action (VLA) Models}

While 2D VLA models have been extensively studied, 3D VLA models
\cite{zhen20243d,yang2025fp3,li2025pointvla,
qu2025spatialvla} remain relatively underexplored.
Zhen \textit{et al.}~\cite{zhen20243d} build 3D-VLA on top of a 3D-based large
language model (LLM) and train it to perform 3D reasoning, multimodal goal
generation, and robot planning.
Lift3D~\cite{jia2024lift3d} enhances 2D foundation models
(\eg, DINOv2~\cite{oquab2023dinov2}) with implicit and explicit 3D robotic
representations for learning 3D manipulation policies.
FP3~\cite{yang2025fp3} employs a transformer to fuse information from point
clouds, proprioceptive states, and language instructions.
PointVLA~\cite{li2025pointvla} uses a VLM and a point-cloud encoder to process
2D images and 3D point clouds, respectively, and injects the resulting 3D
features additively into a few selected blocks of an otherwise frozen action
expert, thereby avoiding retraining of the pre-trained VLA.
SpatialVLA~\cite{qu2025spatialvla} introduces Ego3D positional encoding to
inject 3D information into 2D visual observations and adopts adaptive action
grids to represent robot motions in a more transferable manner.
In contrast to these architectural modifications, \method incorporates 3D
spatial priors without introducing a dedicated 3D encoder or modifying the
core architecture of the VLM: it projects 3D point clouds into multi-view
orthographic images~\cite{goyal2023rvt,goyal2024rvt} and formulates action
prediction as 2D spatial heatmap estimation, thereby keeping both inputs
and outputs within the native 2D domain of the pre-trained VLM.
A concurrent work, OG-VLA~\cite{singh2025ogvla}, explores a similar
orthographic-projection design, albeit generating heatmaps with an
auxiliary image-diffusion decoder.
However, these 3D VLA models operate without any form of memory, which
motivates the line of work on memory-dependent manipulation reviewed next.



\subsection{Memory-Dependent Manipulation}
Most of the above-mentioned VLA models and 3D manipulation policies,
including our \method, adopt a strictly Markovian formulation, predicting
actions solely from the current observation.
Such a formulation becomes inadequate when a task requires temporal
context, or when critical spatial geometry is occluded during execution.
On the one hand, to incorporate \emph{temporal context}, early approaches
attended to the entire observation history~\cite{guhur2023instruction},
resulting in computational costs that scaled poorly with episode length.
Recent methods instead maintain bounded or structured memory through
mechanisms such as explicit memory banks
(\eg, SAM2Act+~\cite{fang2025sam2act}), generative
world models~\cite{li2026lingbotva,yang2026wla,yang2026memorywam},
hierarchical planners~\cite{chen2026rmbench}, visual
traces~\cite{zheng2024tracevla}, or specialized cognitive or gated memory
modules~\cite{lei2025robomemory,shi2026memoryvla,gao2026gatedmemory}.
On the other hand, to address \emph{spatial occlusion}, existing solutions
typically track object poses explicitly over time~\cite{zhou2024yoso} or
maintain persistent geometric representations of the
workspace~\cite{zheng2026memworld}.
In contrast to the above methods, \memmethod offers a simple yet effective
solution to both challenges.
By jointly designing temporal memory for retaining historical context and
spatial memory for recovering occluded geometry, \memmethod achieves strong
performance on memory-dependent manipulation tasks using only a lightweight
attention module.